\begin{textsample}{POS Dim 1 – vem\_ed\_27 – Score 10.00 – vem\_ed\_27/t144.txt}  \label{ex:f1_pos_020}
Coordenador dos PCIs / Cesu de Norte a Sul da América Latina Em novembro e dezembro de 2024, Osvaldo Succi Junior, coordenador dos PCIs / Cesu, realizou visitas técnicas, conferências e workshops nos extremos Norte e Sul da América Latina: México e Chile.
No México, fez a conferência magistral “La Inteligência Artificial em entornos COIL”, no dia 27 de novembro, durante a 2ª Jornada COIL UNAM 2024, cujo tema \textbf{foi} “Aprender por el mundo sin salir de \textbf{casa}”.
A apresentação está disponível a \textbf{partir} do minuto 43 no link da gravação no Facebook: https://www.facebook.com/watch/live/?ef=watch\_permalink\&v=1282050929490379 Além da visita técnica à Universidad Nacional Autónoma de México (UNAM) e da conferência magistral, Succi Junior ofereceu o workshop “Consolidación \textbf{institucional} de um \textbf{programa} COIL”.
A programação incluiu, ainda, a palestra “Una introducción a los proyectos COIL” para um grupo de professores da Faculdade de Estudos Superiores Zaragoza e a apresentação “COIL en la formación universitaria” para docentes da Coordenação de Universidade Aberta e Educação Digital. \textbf{continuação} No Chile, o coordenador dos PCIs / Cesu ministrou um workshop sobre internacionalização para administradores da DUOC UC em 10 de dezembro; também uma palestra sobre Inteligência Artificial Generativa para professores da universidade chilena (no dia 11) e uma apresentação no 1º Seminário Internacional de Learn Chile.
Esse último evento, realizado em 12 de dezembro e transmitido ao vivo pelo YouTube, abordou o tema “Do \textbf{local} ao global: \textbf{conectando} a docência técnica \textbf{profissional} para um mundo globalizado”.
As principais instituições de ensino técnico \textbf{profissional} do Chile participaram do evento: DUOC UC, Inacap, Universidad Técnica Federico Santa María e Santo Tomás.
O objetivo do seminário \textbf{foi} refletir sobre desafios e oportunidades da formação técnica-profissional em \textbf{contextos} internacionais.
A rede Learn Chile \textbf{conta} com cerca de 20 instituições de ensino \textbf{superior}, incluindo universidades e institutos técnico-profissionais, e busca impulsionar a internacionalização da oferta acadêmica chilena e posicionar o país como destino \textbf{educativo} de excelência para estudantes internacionais.
A gravação do evento está disponível em: https://www.youtube.com/live/x6eTux-lIXI

% matched lemmas: casa, conectar, contar, contexto, continuação, educativo, institucional, local, partir, profissional, programa, ser, superior
\end{textsample}
